% Gerado por tcc/analise/resultados_c3.py -- não edite à mão; rode o script de novo.
\begin{table}[htb]
\caption{Endividamento das empresas familiares (F) e não familiares (NF), 2018--2025}\label{tab:completo}
\centering\footnotesize\setlength{\tabcolsep}{4pt}
\begin{tabular}{lrrcccc}
\toprule
Indicador & F & NF & n & Anos F $<$ NF & p (t) & p (MW) \\
\midrule
Dívida bruta / Ativo total & 26,1\% & 30,6\% & 10 $\times$ 6 & 8 de 8 & 0,556 & 0,255 \\
Exigível total / Ativo total & 55,4\% & 56,2\% & 10 $\times$ 6 & 4 de 8 & 0,885 & 0,957 \\
Dívida bruta / (Dívida bruta + PL) & 37,2\% & 43,5\% & 10 $\times$ 6 & 6 de 8 & 0,563 & 0,255 \\
Dívida líquida / EBITDA & 1,4x & 4,4x & 10 $\times$ 6 & 5 de 8 & 0,285 & 0,416 \\
Dívida CP / Dívida total & 20,4\% & 38,2\% & 10 $\times$ 5 & 8 de 8 & 0,112 & 0,058$^{*}$ \\
\bottomrule
\end{tabular}
\fonte{Elaborado pelo autor (2026), com dados da Economática.}
\nota{Testes bicaudais sobre a média de cada empresa nas datas de 31/12 do período. p (MW): teste de Mann-Whitney, aproximação normal com correção de continuidade. $^{***}$, $^{**}$ e $^{*}$: significativo a 1\%, 5\% e 10\%. Valores: média simples das médias por empresa. Na Dívida CP / Dívida total, a Bradsaúde é excluída por não ter dívida financeira; na Dívida líquida / EBITDA, a Biomm tem um único ano com EBITDA positivo (2025).}
\end{table}
